% !TeX program = lualatex
% !TeX encoding = utf8
% !TeX spellcheck = uk_UA
% !BIB program = biber

\documentclass{PySCFBook}


\begin{document}

\pagestyle{main}
%\multiinclude{
%}[]

%%% ============================================================
%\section{Молекулярні властивості як похідні енергії}
%%% ============================================================
%
%\subsection{Основна ідея}
%Молекулярні властивості у квантовій хімії визначаються як аналітичні похідні
%енергії $E$ по зовнішніх параметрах $x$, які характеризують дію поля або
%зміщення системи. Загальна розкладка має вигляд:
%\[
%E(x)
%= E^{(0)} + E^{(1)} x + \tfrac{1}{2} E^{(2)} x^2 + \tfrac{1}{6} E^{(3)} x^3 + \dots
%\]
%Тоді коефіцієнти розкладу — це властивості відповідного порядку:
%\[
%E^{(n)} = \left.\frac{d^nE}{dx^n}\right|_{x=0}.
%\]
%
%\begin{itemize}
%  \item $E^{(1)}$ --- лінійна відповідь системи (дипольний момент, сили);
%  \item $E^{(2)}$ --- квадратична відповідь (поляризованість, гессіан);
%  \item $E^{(3)}$ --- нелінійна відповідь (гіперполяризованість, анеомалії тощо).
%\end{itemize}
%
%\subsection{Загальне співвідношення для будь-якого варіаційного методу}
%Якщо хвильова функція $\Psi$ залежить від зовнішнього параметра $x$
%(наприклад, електричного поля, магнітного поля або координат ядер),
%то енергія
%\[
%E(x) = \langle \Psi(x) | \hat H(x) | \Psi(x) \rangle.
%\]
%
%Її похідні мають вигляд:
%\[
%\frac{dE}{dx}
%= \left\langle \Psi \middle| \frac{\partial \hat H}{\partial x} \middle| \Psi \right\rangle
%+ 2\,\Re \left\langle \frac{\partial \Psi}{\partial x} \middle| \hat H \middle| \Psi \right\rangle.
%\]
%Якщо метод варіаційний (HF, KS-DFT, CASSCF), то
%\(\hat H\Psi=E\Psi\) і $\langle \Psi | \Psi \rangle = 1$, тому
%\[
%\boxed{
%\frac{dE}{dx}
%= \left\langle \Psi \middle| \frac{\partial \hat H}{\partial x} \middle| \Psi \right\rangle.
%}
%\]
%Тобто \textbf{перша похідна енергії дорівнює середньому значенню похідної гамільтоніана}.
%
%\subsection{Приклад 1: Дипольний момент}
%Якщо зовнішній параметр --- однорідне електричне поле $\mathbf F$, то
%гамільтоніан має вигляд
%\[
%\hat H(\mathbf F) = \hat H_0 - \mathbf F \cdot \hat{\boldsymbol{\mu}},
%\]
%де $\hat{\boldsymbol{\mu}} = -\sum_i e\,\mathbf r_i$ --- оператор дипольного моменту.
%
%Похідна енергії за компонентами поля дає середнє значення оператора:
%\[
%\boxed{
%\boldsymbol{\mu}
%= -\frac{\partial E}{\partial \mathbf F}
%= \left\langle \Psi \middle| \hat{\boldsymbol{\mu}} \middle| \Psi \right\rangle.
%}
%\]
%Це --- перший порядок відгуку системи (лінійний ефект).
%
%\subsection{Приклад 2: Ядерні сили і гессіан}
%Якщо зовнішній параметр --- координати ядер $\mathbf R = \{R_A\}$,
%то похідна енергії по $R_A$ визначає силу:
%\[
%\mathbf F_A = -\frac{dE}{d\mathbf R_A}
%= -\left\langle \Psi \middle| \frac{\partial \hat H}{\partial \mathbf R_A} \middle| \Psi \right\rangle
%- \frac{\partial V_{\text{NN}}}{\partial \mathbf R_A}.
%\]
%Це --- \textbf{похідна першого порядку} (градієнт).
%
%Друга похідна по координатах ядер --- \textbf{ядерний гессіан}:
%\[
%\boxed{
%H_{AB} = \frac{d^2E}{dR_A\,dR_B}.
%}
%\]
%Він описує зміну сили на ядрі $A$ при малому зміщенні ядра $B$.
%Саме цей тензор визначає нормальні коливання і частоти:
%\[
%\boldsymbol{\omega}^2 = \frac{\mathbf H}{\mathbf M}.
%\]
%
%\subsection{Фізичний зміст}
%\begin{itemize}
%  \item Перша похідна $dE/dR_A$ --- лінійний відгук енергії на зміщення ядра (сила).
%  \item Друга похідна $d^2E/dR_A dR_B$ --- лінійний відгук \emph{сил} на зміщення іншого ядра (жорсткість системи).
%\end{itemize}
%
%\subsection{Порівняння різних типів похідних}
%\begin{center}
%\begin{tabular}{lcl}
%\toprule
%\textbf{Параметр $x$} & \textbf{Похідна $dE/dx$} & \textbf{Властивість} \\
%\midrule
%Електричне поле $\mathbf F$ & $-\boldsymbol{\mu}$ & Дипольний момент \\
% & $-\frac{d^2E}{d\mathbf F^2}$ & Поляризованість \\
%\hline
%Магнітне поле $\mathbf B$ & $-\boldsymbol{m}$ & Магнітний момент \\
% & $-\frac{d^2E}{d\mathbf B^2}$ & Магнітна сприйнятливість \\
%\hline
%Координати ядер $\mathbf R$ & $-\frac{dE}{d\mathbf R}$ & Сили на ядрах \\
% & $\frac{d^2E}{d\mathbf R^2}$ & Ядерний гессіан \\
%\bottomrule
%\end{tabular}
%\end{center}
%
%\subsection{Ключовий висновок}
%\[
%\boxed{
%\text{Будь-яка молекулярна властивість = похідна енергії по зовнішньому параметру.}
%}
%\]
%Аналітичні методи похідних (CPHF/CPKS) дають змогу обчислювати ці похідні без чисельного диференціювання, через рівняння відгуку хвильової функції.
%
%%% ============================================================
%\section{Аналітичні похідні енергії в теорії Hartree–Fock}
%%% ============================================================
%
%\subsection{Варіаційний принцип і орбітальні параметри}
%У методі Хартрі–Фока хвильова функція системи описується детермінантом Слейтера,
%побудованим із молекулярних орбіталей
%\[
%\phi_i = \sum_{\mu} C_{\mu i}\, \chi_\mu,
%\]
%де $\chi_\mu$ --- атомні базисні функції, а $C_{\mu i}$ --- коефіцієнти розкладу.
%
%Енергія однодетермінантної хвильової функції:
%\[
%E_{\text{HF}}[\mathbf C]
%= \mathrm{Tr}[\mathbf D\,\mathbf h]
%+ \tfrac{1}{2}\mathrm{Tr}[\mathbf D\,\mathbf G(\mathbf D)].
%\]
%
%Тут
%\[
%D_{\mu\nu} = 2 \sum_i C_{\mu i} C_{\nu i},
%\quad
%\mathbf h = \mathbf T + \mathbf V_{\text{ne}},
%\quad
%\mathbf G(\mathbf D) = \mathbf J(\mathbf D) - \tfrac{1}{2}\mathbf K(\mathbf D).
%\]
%
%Варіаційна мінімізація $E_{\text{HF}}[\mathbf C]$ за обмеженням ортонормальності орбіталей
%$\mathbf C^\mathrm{T}\mathbf S\,\mathbf C = \mathbf I$
%призводить до рівнянь Хартрі–Фока:
%\[
%\mathbf F\,\mathbf C = \mathbf S\,\mathbf C\,\boldsymbol{\varepsilon},
%\qquad
%\mathbf F = \mathbf h + \mathbf G(\mathbf D).
%\]
%\[
%\Rightarrow
%\frac{\partial E}{\partial \mathbf C} = 0.
%\]
%
%---
%
%\subsection{Залежність енергії від зовнішнього параметра}
%Нехай система залежить від зовнішнього параметра $x$
%(наприклад, координати ядра $R_A$, електричного поля $F_\alpha$ тощо):
%\[
%E_{\text{HF}} = E_{\text{HF}}(x, \mathbf C(x)).
%\]
%
%Повна похідна енергії:
%\[
%\frac{dE}{dx}
%=
%\frac{\partial E}{\partial x}
%+ \sum_{p}
%\frac{\partial E}{\partial C_p}\frac{dC_p}{dx}.
%\]
%
%Оскільки $\frac{\partial E}{\partial C_p}=0$ у стаціонарній точці (варіаційний принцип),
%перша похідна спрощується:
%\[
%\boxed{
%\frac{dE}{dx} = \frac{\partial E(x,\mathbf C_0)}{\partial x}.
%}
%\]
%Це й дає відомий результат --- градієнт HF обчислюється лише через похідні інтегралів:
%\[
%\frac{dE}{dR_A}
%= \mathrm{Tr}[\mathbf D\,\mathbf h^{(A)}]
%+ \tfrac{1}{2}\mathrm{Tr}[\mathbf D\,\mathbf G^{(A)}]
%- \mathrm{Tr}[\mathbf F\,\mathbf D\,\mathbf S^{(A)}]
%+ \frac{\partial V_{\text{NN}}}{\partial R_A}.
%\]
%
%---
%
%\subsection{Друга похідна: поява відгуку орбіталей}
%Для другої похідної (напр., по $R_A$ і $R_B$):
%\[
%\frac{d^2E}{dR_A\,dR_B}
%=
%\frac{d}{dR_B}
%\left(\frac{\partial E}{\partial R_A}\right)
%=
%\frac{\partial^2E}{\partial R_A\,\partial R_B}
%+ \sum_p
%\frac{\partial^2E}{\partial R_A\,\partial C_p}
%\frac{dC_p}{dR_B}.
%\]
%
%Появляється похідна $\frac{d\mathbf C}{dR_B}$ — це зміна орбіталей при зміщенні ядра.
%Вона визначається з умови мінімуму енергії:
%\[
%\frac{\partial^2E}{\partial \mathbf C\,\partial \mathbf C}\frac{d\mathbf C}{dR_B}
%= -\,\frac{\partial^2E}{\partial R_B\,\partial \mathbf C}.
%\]
%
%Це рівняння — форма \textbf{CPHF (Coupled Perturbed HF)}.
%Позначаючи
%\[
%\mathbf H = \frac{\partial^2E}{\partial \mathbf C\,\partial \mathbf C},
%\quad
%\mathbf g_B = \frac{\partial^2E}{\partial R_B\,\partial \mathbf C},
%\]
%отримуємо:
%\[
%\mathbf H\,\frac{d\mathbf C}{dR_B} = -\,\mathbf g_B,
%\qquad
%\frac{d\mathbf C}{dR_B} = -\,\mathbf H^{-1}\mathbf g_B.
%\]
%
%---
%
%\subsection{Остаточна формула для другої похідної енергії}
%Підставляючи цей результат у вираз для $\frac{d^2E}{dR_A dR_B}$, маємо:
%\[
%\boxed{
%\frac{d^2E}{dR_A\,dR_B}
%=
%\frac{\partial^2E}{\partial R_A\,\partial R_B}
%- \mathbf g_A^\mathrm{T}\,\mathbf H^{-1}\,\mathbf g_B
%+ \frac{\partial^2V_{\text{NN}}}{\partial R_A\,\partial R_B}.
%}
%\]
%
%\begin{itemize}
%  \item Перший доданок --- \emph{заморожені орбіталі}, тобто друга похідна інтегралів по координатам.
%  \item Другий доданок --- \emph{електронний відгук}, знайдений із CPHF.
%  \item Третій --- \emph{ядерно–ядерна взаємодія}.
%\end{itemize}
%
%---
%
%\subsection{Матрична форма рівнянь CPHF}
%У MO-представленні з параметрами ротацій $\kappa_{ai}$ (віртуальні $a$, зайняті $i$)
%CPHF має вигляд:
%\[
%\sum_{bj}
%\left[(\varepsilon_a-\varepsilon_i)\delta_{ab}\delta_{ij}
%+ 2\langle aj||ib\rangle\right]\,U_{bj}^{(R)}
%= -\,2\,F_{ai}^{(R)}.
%\]
%
%Символічно:
%\[
%(\mathbf A+\mathbf B)\,\mathbf U^{(R)} = -\,\mathbf g_R,
%\]
%де \(\mathbf H = 2(\mathbf A+\mathbf B)\) --- електронний гессіан у просторі орбітальних параметрів.
%
%Розв'язок дає \(\mathbf U^{(R)} = -\mathbf H^{-1}\mathbf g_R\),
%що підставляється у формулу другої похідної енергії.
%
%---
%
%\subsection{Підсумок}
%\[
%\boxed{
%\frac{d^2E_{\text{HF}}}{dR_A\,dR_B}
%=
%\mathrm{Tr}[\mathbf D\,\mathbf h^{(AB)}]
%+ \tfrac{1}{2}\mathrm{Tr}[\mathbf D\,\mathbf G^{(AB)}]
%+ \frac{\partial^2V_{\text{NN}}}{\partial R_A\partial R_B}
%- \mathbf g_A^\mathrm{T}\mathbf H^{-1}\mathbf g_B.
%}
%\]
%
%Ця формула поєднує прямі похідні інтегралів і лінійний відгук електронної системи.
%Вона є базовою для всіх аналітичних гессіанів у квантовій хімії.



% !TeX program = lualatex
% !TeX encoding = utf8
% !TeX spellcheck = uk_UA
% !TeX root =../PyscfBook.tex

%=========================================================
%\Opensolutionfile{answer}[\currfilebase/\currfilebase-Answers]
\chapter{Теорія відгуку та властивості молекул}\label{\currfilebase}
%=========================================================

%% ========================================================
\section{Теорія відгуку: загальна концепція}
%% ========================================================

Розглянемо систему, описану гамільтоніаном $\hat{H}_0$, яка перебуває в стаціонарному стані $\Psi_0$ з енергією $E_0$.
Якщо на неї діє зовнішнє збурення $\lambda \hat{V}$, то повний гамільтоніан набуває вигляду
\[
\hat{H}(\lambda) = \hat{H}_0 + \lambda \hat{V}.
\]
Параметр $\lambda$ визначає \emph{силу збурення}: це може бути електричне або магнітне поле, деформація геометрії, зміна спіну тощо.
Відповідно, хвильова функція та енергія системи теж залежать від $\lambda$:
\[
\Psi(\lambda) = \Psi^{(0)} + \lambda\,\Psi^{(1)} + \tfrac{1}{2}\lambda^2\Psi^{(2)} + \cdots,
\]
\[
E(\lambda) = E^{(0)} + \lambda E^{(1)} + \tfrac{1}{2}\lambda^2 E^{(2)} + \cdots.
\]

%% --------------------------------------------------------
\subsection{Розклад енергії за збуренням}
%% --------------------------------------------------------

Підставимо цей розклад у рівняння Шредінґера
\[
\hat{H}(\lambda)\Psi(\lambda) = E(\lambda)\Psi(\lambda),
\]
і прирівняємо члени однакового порядку за $\lambda$.
Для нульового, першого й другого порядків отримаємо:
\[
\begin{aligned}
\hat{H}_0 \Psi^{(0)} &= E^{(0)} \Psi^{(0)}, \\[3pt]
(\hat{H}_0 - E^{(0)})\Psi^{(1)} &= -(\hat{V} - E^{(1)})\Psi^{(0)}, \\[3pt]
(\hat{H}_0 - E^{(0)})\Psi^{(2)} &= -(\hat{V} - E^{(1)})\Psi^{(1)} + E^{(2)}\Psi^{(0)}.
\end{aligned}
\]

А поправки до енергії виражаються через оператор збурення $\hat{V}$ як:
\[
E^{(1)} = \langle \Psi^{(0)} | \hat{V} | \Psi^{(0)} \rangle,
\qquad
E^{(2)} = \langle \Psi^{(0)} | \hat{V} | \Psi^{(1)} \rangle.
\]

Це --- основа \textbf{теорії збурень}. Покажемо, як з неї випливає теорія відгуку.

Оскільки енергію системи можна виразити як сподіване значення гамільтоніана системи:

\[
E(\lambda) = \frac{\langle \Psi(\lambda) | \hat{H}(\lambda) | \Psi(\lambda) \rangle}
{\langle \Psi(\lambda) | \Psi(\lambda) \rangle}.
\]

То згідно теореми \textbf{Гельмана-Фейнмана}, похідна енергії по параметру збурення $\lambda$ має вигляд:
\[
\frac{dE}{d\lambda}
= \bracket<\Psi(\lambda) | \frac{d\hat{H}}{d\lambda} | \Psi(\lambda) >.
\]


Якщо далі підставити $\hat{H}(\lambda) = \hat{H}_0 + \lambda \hat{V}$, то отримаємо:
\[
\left.\frac{dE}{d\lambda}\right|_{\lambda=0}
= \langle \Psi^{(0)} | \hat{V} | \Psi^{(0)} \rangle
\;\equiv\; E^{(1)}.
\]
Отже, $E^{(1)}$ --- це не просто коефіцієнт при $\lambda$, а \emph{справжня похідна енергії за збуренням}.
Аналогічно, друга похідна
\[
\frac{d^2E}{d\lambda^2} = 2\langle \Psi^{(0)} | \hat{V} | \Psi^{(1)} \rangle.
\]


Покажемо далі, що \emph{усі властивості молекули можна виразити через похідні енергії за параметром, який описує збурення системи.}

%% --------------------------------------------------------
\subsection{Від енергії до властивостей}
%% --------------------------------------------------------

Нехай параметр~$F$ позначає будь-яку \emph{зовнішню дію на систему}:
це може бути електричне або магнітне поле, зміна координат ядер, деформація потенціалу чи спінове збурення.
Тоді повна енергія молекули є аналітичною функцією цього параметра:
\[
E(\mathbf{F}) = E_0
+ \left( \frac{\partial E}{\partial F_\alpha} \right) F_\alpha
+ \frac{1}{2} \left( \frac{\partial^2 E}{\partial F_\alpha \partial F_\beta} \right) F_\alpha F_\beta
+ \cdots
\]
де $F_\alpha$ --- компоненти зовнішньої дії  на систему.

Відповідні похідні визначають фізичні властивості:
\[
\begin{aligned}
\mu_\alpha &= - \frac{\partial E}{\partial E_\alpha}
\quad &&\text{(дипольний момент)},\\[3pt]
\alpha_{\alpha\beta} &= - \frac{\partial^2 E}{\partial E_\alpha \partial E_\beta}
\quad &&\text{(поляризованість)},\\[3pt]
\sigma_{\alpha\beta} &= - \frac{\partial^2 E}{\partial B_\alpha \partial M_\beta}
\quad &&\text{(магнітне екранування)},\\[3pt]
F_{A\alpha} &= - \frac{\partial E}{\partial R_{A\alpha}}
\quad &&\text{(сила на ядрі~$A$)},\\[3pt]
H_{AB}^{\alpha\beta} &= \frac{\partial^2 E}{\partial R_{A\alpha}\partial R_{B\beta}}
\quad &&\text{(гессіан, що визначає частоти коливань)}.
\end{aligned}
\]

Таким чином, як електричні, магнітні, так і механічні властивості молекули
описуються єдиним математичним підходом — через похідні енергії за параметрами, що збурюють систему.
Власне, \textbf{теорія відгуку} --- це узагальнення теорії збурень, в якій
замість формального розкладу хвильової функції розглядається її \emph{варіаційний відгук} на зовнішнє поле.


\begin{table}[h!]
\centering
\caption{Зв’язок порядку похідних енергії з молекулярними властивостями. Тут $q$ --- координати, $E$ --- зовнішнє електричне поле, $B$ --- зовнішнє магнітне поле, $I$ --- магнітне поле ядра.}\label{tab:prop}
\begin{tblr}{
%  cells = {font=\scriptsize},
  colspec={Q[c,m ]Q[c,m] Q[c,m] Q[c,m] Q[l,m]},
  cell{1}{1} = {c=4}{c},
  cell{1}{5} = {}{c,m},
  cell{2-Z}{1-4}={}{fg=blue, cmd=\bfseries},
  stretch = -1,
  hline{2} = {2pt},
  hline{3} = {1pt},
  hline{Z} = {1pt},
}
$n$-та похідна енергії &       &       &       & {Збурення} \\
$n_R$ & $n_E$ & $n_B$ & $n_I$ & Властивість \\
1 & 0 & 0 & 0 & Сила \\
2 & 0 & 0 & 0 & Частоти гармонічних коливань \\
3 & 0 & 0 & 0 & Ангармонічні поправки \\
0 & 1 & 0 & 0 & Електричний дипольний момент $\vec{\mu}_e$ \\
0 & 2 & 0 & 0 & Електрична поляризованість $\chi_e$ \\
0 & 3 & 0 & 0 & Електрична надполяризованість \\
0 & 0 & 1 & 0 & Магнітна сприйнятливість $\chi_m$ \\
0 & 0 & 2 & 1 & Константа надтонкого зв’язку (\texttt{EPR}) \\
0 & 0 & 0 & 2 & Стала спін-спінової взаємодії ядер \\
1 & 1 & 0 & 0 & Інтенсивності ІЧ-спектрів \\
2 & 1 & 0 & 0 & Інтенсивності обертонів в ІЧ-спектрах \\
1 & 2 & 0 & 0 & Інтенсивності Раманівських спектрів \\
2 & 2 & 0 & 0 & Інтенсивності обертонів у Раманівських спектрах \\
0 & 1 & 1 & 0 & Циркулярний дихроїзм \\
0 & 2 & 1 & 0 & Магнітний циркулярний дихроїзм \\
0 & 0 & 1 & 1 & Хімічний зсув (\texttt{NMR}) \\
\end{tblr}
\end{table}

Концепція теорії відгуку надзвичайно універсальна: практично всі молекулярні властивості можна розглядати як похідні енергії за зовнішніми параметрами --- електричним або магнітним полем, координатами ядер тощо.
Нижче наведено узагальнену схему, що показує, як похідні різних порядків визначають ті чи інші властивості (див. табл \ref{tab:prop}).

Як видно, усі ці властивості --- від сил та частот до магнітного екранування ---
описуються похідними енергії різних порядків за різними збуреннями.
Таким чином, \emph{усі властивості молекули є лише різними проявами одного й того ж принципу: відгуку хвильової функції на зовнішнє збурення.}

%% --------------------------------------------------------
\subsection{Формалізм варіацій орбіталей}
%% --------------------------------------------------------

У загальній теорії відгуку рівняння Шредінґера описує зміну повної хвильової функції~$\Psi(\lambda)$ під дією збурення.
Однак у багаточастинкових системах (молекулах) така функція залежить від великої кількості координат електронів,
і безпосередньо працювати з нею практично неможливо.

У методі Хартрі–Фока ця складна залежність замінюється набором одноелектронних функцій --- молекулярних орбіталей~$\phi_i$,
з яких будується детермінант Слейтера.
Саме ці орбіталі і є основними варіаційними змінними, що визначають енергію системи.

Тому, щоб реалізувати теорію відгуку на практиці, необхідно перейти від рівняння Шредінґера до орбітального формалізму,
у якому збурення описується через варіації орбіталей, а зміна енергії --- через відповідні зміни операторів Фока.

У методі Хартрі-Фока хвильова функція є детермінантом із молекулярних орбіталей~$\phi_i$, що задовольняють рівняння
\[
\hat{F}\,\phi_i = \varepsilon_i \phi_i.
\]
При прикладенні зовнішнього поля $\lambda \hat{V}$ орбіталі змінюються:
\[
\phi_i \;\to\; \phi_i + \lambda\,\phi_i^{(1)} + \mathcal{O}(\lambda^2),
\]
де $\phi_i^{(1)}$ --- \emph{варіації орбіталей}, або \textbf{збурений відгук} системи.
Підставивши це у рівняння Хартрі---Фока й розклавши за степенями $\lambda$, отримаємо лінійні рівняння для відгуку:
\[
(\hat{F}-\varepsilon_i)\phi_i^{(1)} = -\hat{F}^{(1)}\phi_i,
\]
які відомі як \textbf{рівняння CPHF} (Coupled-Perturbed Hartree--Fock).

Розв’язання цих рівнянь дозволяє визначити, як змінюються орбіталі під дією зовнішнього поля,
а отже --- обчислити похідні енергії, тобто молекулярні властивості.
У методах на основі функціонала густини (\texttt{DFT}) аналогічні рівняння мають назву \textbf{CPKS} (Coupled–Perturbed Kohn--Sham).

На практиці цей формалізм реалізовано в пакеті~\texttt{PySCF},
де користувач може вибрати, чи враховувати релаксацію орбіталей у відгуку:
\begin{itemize}
    \item \inlinecode{cphf=False} --- обчислення без урахування відгуку орбіталей (\emph{unrelaxed} властивості);
    система не адаптується до збурення.
    \item \inlinecode{cphf=True} --- повне самозгоджене врахування відгуку (\emph{relaxed} властивості);
    орбіталі релаксують під дією поля, що забезпечує фізично коректні похідні енергії.
\end{itemize}

\input{\currfiledir ElectricProperties.tex}





\end{document}
